The \wikilink{s40--b008}{definition for one element} does not require
the whole extension to be finite. It gives the following general definition.

\begin{definition}\label{separable-algebraic-extension}
An algebraic extension $K/k$ is \emph{separable} if every element of
$K$ is separable over $k$, that is, its minimal polynomial has no
repeated roots in an algebraic closure.
\end{definition}
\begin{proposition}\label{separable-finite-subextensions}
An algebraic extension is separable if and only if each finite
intermediate extension is separable. For a finite extension this
definition agrees with the embedding-count definition in
\wikilink{s40--b012}{the source notes}.
\end{proposition}
\begin{proof}\label{separable-finite-subextensions-proof}
Suppose all elements are separable. A finite intermediate extension
is generated by finitely many such elements. In a tower adjoining
them successively, each new minimal polynomial divides the separable
minimal polynomial over $k$, and is therefore separable.
The simple-extension embedding criterion and the finite tower
criterion in \wikilink{s40--b010}{simple separability} and
\wikilink{s40--b019}{separable towers} show that the resulting finite
extension has the full number of embeddings.
Conversely, if a finite extension has that property, the same tower
criterion applied to $K/k(\alpha)/k$ makes every $k(\alpha)/k$
separable. Its simple-extension criterion makes $\alpha$ separable.
Finally, every element of an algebraic extension lies in the finite
intermediate field $k(\alpha)$, proving the general equivalence.
\end{proof}
\begin{remark}\label{separability-characteristic-zero}
In characteristic zero every algebraic extension is separable, even
an infinite one: apply \wikilink{s43--b001}{the characteristic-zero
irreducibility argument} to each element's finite minimal polynomial.
Embedding counts and finite degree formulas should not be silently
replaced by arithmetic on infinite cardinals.
\end{remark}
